%% Generated by Sphinx.
\def\sphinxdocclass{report}
\documentclass[letterpaper,10pt,english]{sphinxmanual}
\ifdefined\pdfpxdimen
   \let\sphinxpxdimen\pdfpxdimen\else\newdimen\sphinxpxdimen
\fi \sphinxpxdimen=.75bp\relax
\ifdefined\pdfimageresolution
    \pdfimageresolution= \numexpr \dimexpr1in\relax/\sphinxpxdimen\relax
\fi
\newdimen\sphinxremdimen\sphinxremdimen = 10pt
%% let collapsible pdf bookmarks panel have high depth per default
\PassOptionsToPackage{bookmarksdepth=5}{hyperref}

\PassOptionsToPackage{booktabs}{sphinx}
\PassOptionsToPackage{colorrows}{sphinx}

\PassOptionsToPackage{warn}{textcomp}
\usepackage[utf8]{inputenc}
\ifdefined\DeclareUnicodeCharacter
% support both utf8 and utf8x syntaxes
  \ifdefined\DeclareUnicodeCharacterAsOptional
    \def\sphinxDUC#1{\DeclareUnicodeCharacter{"#1}}
  \else
    \let\sphinxDUC\DeclareUnicodeCharacter
  \fi
  \sphinxDUC{00A0}{\nobreakspace}
  \sphinxDUC{2500}{\sphinxunichar{2500}}
  \sphinxDUC{2502}{\sphinxunichar{2502}}
  \sphinxDUC{2514}{\sphinxunichar{2514}}
  \sphinxDUC{251C}{\sphinxunichar{251C}}
  \sphinxDUC{2572}{\textbackslash}
\fi
\usepackage{cmap}
\usepackage[T1]{fontenc}
\usepackage{amsmath,amssymb,amstext}
\usepackage{babel}



\usepackage{tgtermes}
\usepackage{tgheros}
\renewcommand{\ttdefault}{txtt}



\usepackage[Glenn]{fncychap}
\usepackage{sphinx}

\fvset{fontsize=auto}
\usepackage{geometry}


% Include hyperref last.
\usepackage{hyperref}
% Fix anchor placement for figures with captions.
\usepackage{hypcap}% it must be loaded after hyperref.
% Set up styles of URL: it should be placed after hyperref.
\urlstyle{same}

\addto\captionsenglish{\renewcommand{\contentsname}{Contents:}}

\usepackage{sphinxmessages}
\setcounter{tocdepth}{2}



\title{Grokking Algorithms Concepts}
\date{Oct 06, 2026}
\release{1.0.0}
\author{Ali Izadi}
\newcommand{\sphinxlogo}{\vbox{}}
\renewcommand{\releasename}{Release}
\makeindex
\begin{document}

\ifdefined\shorthandoff
  \ifnum\catcode`\=\string=\active\shorthandoff{=}\fi
  \ifnum\catcode`\"=\active\shorthandoff{"}\fi
\fi

\pagestyle{empty}
\sphinxmaketitle
\pagestyle{plain}
\sphinxtableofcontents
\pagestyle{normal}
\phantomsection\label{\detokenize{index::doc}}



\chapter{Problems}
\label{\detokenize{index:problems}}

\section{Problem 1:}
\label{\detokenize{index:problem-1}}

\subsection{I can’t figure out the binary search implementation!?}
\label{\detokenize{index:i-can-t-figure-out-the-binary-search-implementation}}
\begin{sphinxadmonition}{note}{Try to solve:}

\sphinxAtStartPar
At first I try to implement binary search without AI help or viewing book’s implementation.
But this is not worked. After that I use AI to get a correct implementation but my code output
was not correct! Finally I used book’s implementation in page 9 to correct my different program.
\end{sphinxadmonition}


\section{Problem 2:}
\label{\detokenize{index:problem-2}}

\subsection{I was a little confused about logarithms!}
\label{\detokenize{index:i-was-a-little-confused-about-logarithms}}
\begin{sphinxadmonition}{note}{Try to solve:}
\begin{enumerate}
\sphinxsetlistlabels{\arabic}{enumi}{enumii}{}{.}%
\item {} 
\sphinxAtStartPar
Read the book’s definition for logarithms.

\item {} 
\sphinxAtStartPar
Watch some tutorials in KhanAcademy about logarithms.

\item {} 
\sphinxAtStartPar
Solve some practices in KhanAcademy about logarithms.

\end{enumerate}
\end{sphinxadmonition}


\section{Problem 3:}
\label{\detokenize{index:problem-3}}

\subsection{I had no idea to understand the concept of permuatation!}
\label{\detokenize{index:i-had-no-idea-to-understand-the-concept-of-permuatation}}
\begin{sphinxadmonition}{note}{Try to solve:}
\begin{enumerate}
\sphinxsetlistlabels{\arabic}{enumi}{enumii}{}{.}%
\item {} 
\sphinxAtStartPar
Watch one video about permutation in KhanAcademy

\item {} 
\sphinxAtStartPar
Compare this concept to factorial concept (n!) and the ambiguity was resolved.

\end{enumerate}
\end{sphinxadmonition}



\renewcommand{\indexname}{Index}
\printindex
\end{document}